% MA 214 Discussion 12 -- covers Lectures 17 and 18.
%
% Student handout. Page 1 is the concept review; the question set runs from page 2.
% Compile from inside this folder:  pdflatex Discussion12.tex

\documentclass[11pt]{article}
\usepackage[margin=1in]{geometry}
\usepackage{parskip}
\usepackage{fancyhdr}
\usepackage{array}
\usepackage{booktabs}
\usepackage{enumitem}
\usepackage{amsmath}
\usepackage{tabularx}
\usepackage{listings}

\pagestyle{fancy}
\fancyhf{}
\cfoot{\thepage}
\rhead{Discussion 12}
\lhead{MA 214: Applied Statistics}
\renewcommand{\headrulewidth}{.4mm}

\newcommand{\blankline}[1]{\underline{\hspace{#1}}}
\newcolumntype{L}{>{\raggedright\arraybackslash}X}

\lstset{
  basicstyle=\ttfamily\small,
  columns=fullflexible,
  frame=single,
  framerule=0.4pt,
  breaklines=true,
  showstringspaces=false,
  aboveskip=8pt,
  belowskip=8pt
}

% Section header: bold title over a hairline rule.
\newcommand{\parttitle}[1]{%
  \par\medskip\noindent
  {\large\bfseries #1}\par
  \vspace{-.5em}\noindent\rule{\linewidth}{0.4pt}\par\smallskip}

% Writing space.
\newcommand{\answerspace}[1]{\vspace{#1}}

% Flush-left question list: the label runs in at the margin and the body wraps
% back to the margin, so nothing is pushed far to the right.
\newlist{questions}{enumerate}{1}
\setlist[questions]{label=\textbf{Question \arabic*.}, wide=0pt, itemsep=1.4em, topsep=.8em}

\newlist{parts}{enumerate}{1}
\setlist[parts]{label=\textbf{(\Alph*)}, leftmargin=1.6em, labelsep=.45em,
  itemsep=.7em, topsep=.5em}

\begin{document}

\noindent\textbf{Name:}\blankline{2.3in}\hfill\textbf{BUID:}\blankline{1.6in}

\begin{center}
  {\Large\bfseries Discussion 12: From a Grid to a Curve}
\end{center}

\noindent\textbf{Goals:} \textbf{(1)} apply Bayes' rule to events and explain why a base rate can
dominate a positive test, \textbf{(2)} update sequentially and see that order does not matter,
\textbf{(3)} normalize a discrete posterior by hand, \textbf{(4)} use the conjugate update rule and
read the posterior mean as a weighted average, \textbf{(5)} say what a prior is worth in
observations.

\parttitle{Part 1: Concept Review}

{\small

\noindent Everything in Part 2 can be answered from this page. Keep it in front of you.

\begin{enumerate}[label=\textbf{\arabic*.}, wide=0pt, itemsep=.4em, topsep=.5em]

\item \textbf{What is random, and Bayes' rule for events.} Frequentist: the parameter is fixed, the
data random. Bayesian: the data are what you have and the \emph{parameter} gets a distribution, a
\textbf{prior} before and a \textbf{posterior} after. For a hypothesis $H$ and evidence $E$,
\[
P(H \mid E) = \frac{P(E \mid H)\,P(H)}{P(E)},
\qquad
P(E) = P(E \mid H)P(H) + P(E \mid H^c)P(H^c).
\]
When $P(H)$ is small, one positive result often leaves $P(H\mid E)$ small too: the false positives
from the large $H^c$ group swamp the true positives. The prior does the work.

\item \textbf{Sequential updating.} Feed each posterior in as the next prior. The final answer does
not depend on the order the observations arrive in, because the product of prior and likelihoods is
the same either way.

\item \textbf{Bayes' rule for a parameter.} Over a discrete set of candidate values,
$f(\theta \mid x) \propto f(\theta)L(\theta;x)$: multiply prior by likelihood, total the products,
divide through by that total. The denominator is only a normalizing constant, which is why the
proportionality is enough to identify the posterior.

\item \textbf{The three conjugate families.} A \textbf{conjugate} prior makes the posterior come
from the same family, so no normalizing is needed: you just update the parameters.

\begin{center}
\renewcommand{\arraystretch}{1.15}
\begin{tabular}{llll}
\toprule
\textbf{Data} & \textbf{Prior} & \textbf{Posterior} & \textbf{Prior is worth} \\
\midrule
$y$ of $n$ successes & Beta$(a,b)$ & Beta$(a+y,\; b+n-y)$ & $a+b$ trials \\
$n$ obs, $\sigma$ known & N$(\theta_0, \tau^2)$ & N(weighted mean, $1/\text{precision}$) & $\sigma^2/\tau^2$ obs \\
counts over $t$ units & Gamma$(s,r)$ & Gamma$(s+\sum y,\; r+t)$ & $r$ units \\
\bottomrule
\end{tabular}
\end{center}

In the Normal--Normal case precision adds, $1/\tau^2 + n/\sigma^2$, and the posterior mean is
$\big(\theta_0/\tau^2 + n\bar{x}/\sigma^2\big)$ divided by that precision.

\item \textbf{The posterior mean is a weighted average.} In all three,
$\text{posterior mean} = (w \cdot \text{prior mean} + n \cdot \text{data estimate})/(w + n)$, where
the prior weight $w$ is $a+b$, $\sigma^2/\tau^2$, or $r$, and the data weight is $n$ or $t$. So the
posterior always lands \emph{between} the prior mean and the data, nearer whichever carries more
weight, and as $n$ grows different priors converge.

\item \textbf{Four traps.} Reporting prior $\times$ likelihood without dividing by the total.
Reusing the original prior at every step of a sequential update instead of the latest posterior.
Swapping the counts in the Beta update. Saying a posterior probability is a $p$-value.

\end{enumerate}

}

\newpage

\parttitle{Part 2: Question Set}

\begin{questions}

%%%%%%%%%%%%%%%%%%%%%%%%%%%%%%%%%%%%%%
\item \textbf{A detector, and a base rate.} A field station runs an automated detector for the call
of a rare frog. On any given night the species is actually present at a monitoring site with
probability $0.03$. The detector has sensitivity $0.92$ and specificity $0.88$.

\begin{parts}
\item Compute $P(\text{fires})$ by the law of total probability, then
$P(\text{frog} \mid \text{fires})$.
\answerspace{5em}

\item The detector is good, yet your answer to (A) is under $0.20$. Explain where the probability
went, in terms of the two groups of sites.
\answerspace{3.2em}

\item The detector fires on a \emph{second} independent night at the same site. Using your answer to
(A) as the new prior, compute the updated probability. A third firing gives $0.933$; what is the
pattern across the three numbers?
\answerspace{3.5em}

\item Suppose instead the detector fires once and then fails to fire once. Someone claims the answer
depends on which happened first. Explain why it does not.
\answerspace{3.2em}
\end{parts}

%%%%%%%%%%%%%%%%%%%%%%%%%%%%%%%%%%%%%%
\item \textbf{A grid, normalized by hand.} Now the unknown is a number: let $\theta$ be the
probability that a randomly chosen site holds the frog. Restrict $\theta$ to four values. You survey
$10$ sites and find the frog at $3$ of them.

\begin{parts}
\item Complete the table. The likelihoods are given.

\begin{center}
\renewcommand{\arraystretch}{1.4}
\begin{tabular}{ccccc}
\toprule
$\theta$ & Prior $f(\theta)$ & $L(\theta; x=3)$ & Product & Posterior $f(\theta \mid x)$ \\
\midrule
0.10 & 0.30 & 0.0574 & \blankline{1in} & \blankline{1in} \\
0.20 & 0.40 & 0.2013 & \blankline{1in} & \blankline{1in} \\
0.30 & 0.20 & 0.2668 & \blankline{1in} & \blankline{1in} \\
0.40 & 0.10 & 0.2150 & \blankline{1in} & \blankline{1in} \\
\midrule
 & & Total: & \blankline{1in} & 1.000 \\
\bottomrule
\end{tabular}
\end{center}

\item The products total $0.1726$, not $1$. Say what that total is called and why dividing by it is
required.
\answerspace{3.5em}

\item Compute the prior mean and the posterior mean of $\theta$. The data on its own suggest
$3/10 = 0.30$. Where does the posterior mean sit relative to $0.21$ and $0.30$, and why?
\answerspace{3.5em}

\item You only allowed four values of $\theta$, but $\theta$ could be anything in $[0,1]$. What
does Lecture 18 replace the four rows with, and what happens to the ``Total'' column?
\answerspace{3.2em}
\end{parts}

%%%%%%%%%%%%%%%%%%%%%%%%%%%%%%%%%%%%%%
\item \textbf{The same rule, three times.} Each row is a different conjugate model. Fill in the
posterior and the weights.

\begin{center}
\renewcommand{\arraystretch}{1.3}
\begin{tabularx}{\linewidth}{LLL}
\toprule
\textbf{Setting} & \textbf{Posterior} & \textbf{Prior / data weight} \\
\midrule
Beta$(2,8)$; $5$ of $20$ sites & & \\
N$(3000,200^2)$; $n=25$, $\bar{x}=2850$, $\sigma=400$ & & \\
Gamma$(4,2)$; $18$ calls in $6$ nights & & \\
\bottomrule
\end{tabularx}
\end{center}

\begin{parts}
\item Give the posterior \emph{mean} in each of the three rows.
\answerspace{3.2em}

\item For each row, state what the prior is worth in observations, and check it against the weights
you wrote. In the Normal--Normal row that worth is only $4$ observations against $n = 25$, which is
why the posterior mean $2870.7$ sits much nearer $2850$ than $3000$.
\answerspace{3.5em}

\item What single feature makes all three of these ``the same rule''?
\answerspace{3em}
\end{parts}

%%%%%%%%%%%%%%%%%%%%%%%%%%%%%%%%%%%%%%
\item \textbf{Five people, one data set.} Five researchers hold different Beta priors for $\theta$.
All see the same data: the frog at $5$ of $20$ sites, a proportion of $0.25$.

\begin{center}
\renewcommand{\arraystretch}{1.3}
\begin{tabular}{lccc}
\toprule
\textbf{Prior} & \textbf{Prior mean} & \textbf{Post.\ mean, $n=20$} & \textbf{Post.\ mean, $n=200$} \\
\midrule
Beta$(2,8)$, skeptical & 0.200 & 0.2333 & 0.2476 \\
Beta$(8,2)$, optimistic & 0.800 & 0.4333 & 0.2762 \\
Beta$(20,80)$, firm skeptic & 0.200 & 0.2083 & 0.2333 \\
\bottomrule
\end{tabular}
\end{center}

\noindent The $n = 200$ column uses $50$ of $200$ sites, the same proportion $0.25$.

\begin{parts}
\item At $n = 20$ the posterior means run from $0.208$ to $0.433$, a spread of $0.225$. At
$n = 200$ they run from $0.233$ to $0.276$, a spread of $0.043$. State the general principle.
\answerspace{3.2em}

\item Beta$(2,8)$ and Beta$(20,80)$ have the \emph{same} prior mean of $0.200$ but reach different
posteriors. Explain why, using what each prior is worth. Then say whether the optimistic
researcher's $0.433$, well above the observed $0.25$, is an arithmetic mistake, a method mistake,
or neither.
\answerspace{5em}
\end{parts}

%%%%%%%%%%%%%%%%%%%%%%%%%%%%%%%%%%%%%%
\item \textbf{Find the bugs.} The code below is meant to reproduce Question 2, then the Beta update
from Question 3, then two sequential detector firings. It runs, prints plausible numbers, and is
wrong in \textbf{three} separate ways.

\begin{lstlisting}[language=R]
lik <- dbinom(3, 10, theta)
posterior <- prior * lik
posterior

a <- 2; b <- 8; y <- 5; n <- 20
c(a + n, b + y)

p <- 0.03
for (i in 1:2) post <- 0.92*p / (0.92*p + 0.12*(1-p))
post
\end{lstlisting}

\begin{parts}
\item Bug 1 is the \texttt{posterior} line: what do those four numbers sum to, and what is missing?
Bug 2 is \texttt{c(a + n, b + y)}: give the numbers it prints and the numbers it should print.
\answerspace{4.5em}

\item Bug 3 is in the loop. What does \texttt{post} equal after two iterations, and what should it
equal? Name the line that has to change.
\answerspace{3.5em}

\item Which of the three bugs would still be invisible if the loop ran only \emph{one} iteration?
\answerspace{3em}
\end{parts}

\end{questions}

\end{document}
